% =============================================================
%  Chapter 2 — Background Review : SKELETON / SCAFFOLD
%  Status: collaborative outline, not draft prose.
%  Each section lists: Role / POV — Key arguments — References
%  from the project .bib — Consensus candidates worth adding —
%  Watch-fors. References use BibTeX keys from references.bib.
% =============================================================

\chapter{Background Review}
\label{ch:background}

% -------------------------------------------------------------
%  Chapter introduction
% -------------------------------------------------------------
\section*{Chapter Introduction}
\label{sec:ch2-intro}

\textbf{Role.} Tell the reader what the chapter does and how it
maps to the dialectics of Chapter~1.

\textbf{POV.} Three converging literatures — emergence, embodiment,
existing soft/modular systems — define the conceptual neighborhood
for the thesis. Each section narrows scope; \S2.3 closes with the
gap statement that motivates Chapter~3.

\textbf{Points to make (5--7 sentences).}
\begin{itemize}
    \item Map: \S2.1 picks up \emph{Bottom-up vs.\ Top-down}; \S2.2
          picks up \emph{Soft vs.\ Rigid}; \S2.3 surveys robots that
          sit somewhere in those tensions.
    \item Briefly acknowledge the early second-order cybernetics
          and HRI thread (\cite{glanvilleSECONDORDERCYBERNETICS},
          \cite{sheridanHumanRobotInteractionStatus2016}) as
          intellectual scaffolding now reserved for the reflection
          chapter, not the experimental work.
    \item Signal the chapter arc: \emph{control paradigm} $\rightarrow$
          \emph{embodiment principle} $\rightarrow$ \emph{existing systems}
          $\rightarrow$ \emph{gap}.
\end{itemize}

% =============================================================
\section{The Elegance of Emergence}
\label{sec:elegance-emergence}
% =============================================================

\textbf{Role.} Establish the bottom-up control paradigm — that
complex global behaviour can arise from simple local rules — as
a precise computational practice, not a metaphor. Grounds the
\emph{Bottom-up vs.\ Top-down} dialectic from \S1.1 and motivates
the CA-based control direction taken later in the thesis.

\textbf{POV.} Emergence is interesting to roboticists not because it
is mysterious but because it is \emph{cheap} — global behaviour for
free, given the right substrate and rule set. The chapter takes
this engineering posture, not a philosophical one.

% -----------------------------------------
\subsection{Defining emergence: a precise relationship, not a metaphor}
\label{subsec:defining-emergence}

\textbf{Key arguments.}
\begin{itemize}
    \item Emergence is a \emph{relational} property of a system: the
          global state cannot be obtained by a compositional sum of
          agent states, even though it is fully determined by them.
    \item Distinguish emergence from non-linearity in the colloquial
          sense; reach instead for an information-theoretic reading.
    \item Brief tour of the language: self-organization, complexity,
          computational irreducibility.
\end{itemize}

\textbf{References (from your .bib).}
\cite{prokopenkoInformationTheoreticPrimerComplexity2006},
\cite{dehaanHowEmergenceArises2006},
\cite{ladymanWhatComplexSystem2013},
\cite{mitchellComplexityGuidedTour2011},
\cite{bonabeauNaturalArtiCial1999}.

\textbf{Consensus candidates worth adding.} None essential — your
bib is strong here. Anderson 1972 ``More is different'' is the
classic foundational citation if you want the canonical opener.

% -----------------------------------------
\subsection{Computational practices that operationalize emergence}
\label{subsec:computational-practices}

\textbf{Key arguments.}
\begin{itemize}
    \item \textbf{Cellular automata} as the cleanest demonstration:
          discrete lattice + local update rule $\rightarrow$ global
          pattern; minimal substrate showing how much complexity
          local rules can support.
    \item \textbf{Self-organization in biology and physics} — Turing
          morphogenesis, ant colony optimization, swarm intelligence
          — same logic, different substrate.
    \item \textbf{Engineered swarms and programmable matter} —
          emergence moves from simulation into hardware.
\end{itemize}

\textbf{References (from your .bib).}
Swarms / engineered emergence: \cite{rubensteinProgrammableSelfassemblyThousandrobot2014},
\cite{liParticleRoboticsBased2019},
\cite{cademartiriProgrammableSelfassembly2015},
\cite{bonabeauNaturalArtiCial1999},
\cite{camazineSelfOrganizationBiologicalSystems2001}.

Behaviour discovery / characterization in swarms:
\cite{mattsonDiscoveryDeploymentEmergent2025},
\cite{mattsonLeveragingHumanFeedback2023},
\cite{rajbhandariLearningNEATEmergent2024},
\cite{walkerCharacterizingHumanPerception2016},
\cite{IndirectSwarmControl}.

\textbf{Consensus candidates worth adding.}
\begin{itemize}
    \item \emph{Werfel et al.\ 2014, Science} --- termite-inspired
          construction swarms; addresses the inverse problem
          (specify the structure, synthesize the rules) that the
          thesis implicitly inverts back. Foundational and frequently
          cited.
    \item \emph{Carrillo-Zapata et al.\ 2018, Science Robotics} ---
          ``Morphogenesis in robot swarms''; closes the loop between
          biological self-organization and physical robot collectives.
    \item \emph{Bojinov et al.\ 2000, ICRA} --- ``Emergent structures
          in modular self-reconfigurable robots''; pre-dates
          \cite{kubicaComplexBehaviorsLocal2001} which you already
          cite, useful as the historical anchor.
\end{itemize}

% -----------------------------------------
\subsection{From rule sets to embodied substrates}
\label{subsec:rules-to-embodied}

\textbf{Key arguments.}
\begin{itemize}
    \item \emph{Where emergence has stalled in robotics:} most
          demonstrations operate on rigid agents whose only coupling
          to one another is informational (radio, sensing). The
          mechanical and material coupling between bodies is left out.
    \item \emph{What we want next:} emergence operating on physically
          coupled bodies, where ``communication'' between modules is
          partly mechanical and material, not only signalled.
    \item \emph{Bridging work that already exists} — voxel-based
          soft robot CAs, modular soft robots controlled by
          NCA-like distributed rules — but it is still nascent.
\end{itemize}

\textbf{References (from your .bib).}
\cite{romanishinMBlocksThreeDimensional2018},
\cite{wangRoboMatterReconfigurableMultifunctional2024},
\cite{eguchiFluidicSwarmEmbodimentSwarm2025},
\cite{mousaMultifunctionalFluidicUnits},
\cite{chattyEmergentComplexBehaviors2011},
\cite{kubicaComplexBehaviorsLocal2001},
\cite{zhangBioInspiredDistributedNeural}.

\textbf{Consensus candidates worth adding.}
\begin{itemize}
    \item \emph{Cheney et al.\ 2016} --- ``Topological evolution for
          embodied cellular automata''; the title alone earns it the
          slot, and it argues the same body+CA framing you use.
    \item \emph{Nadizar et al.\ 2022} --- ``Collective control of
          modular soft robots via embodied Spiking NCA''; cleanest
          modern bridge between CAs and soft modular bodies.
    \item \emph{Oliveri et al.\ 2021, PNAS} --- ``Continuous learning
          of emergent behavior in robotic matter''; physical-coupling-
          driven learning without communication. Strong support for
          the ``mechanical coupling as a substrate for emergence''
          claim.
\end{itemize}

\textbf{Closing transition.} Emergence shows that complex
behaviour need not be specified from the top. But every example
above treats agents as mechanically independent units exchanging
information. The next section examines a related but distinct claim:
that even within a single body, intelligence can be distributed
across material and form rather than concentrated in a controller.

\textbf{Watch out for.}
\begin{itemize}
    \item Don't get pulled into Wolfram metaphysics.
    \item Don't promise more than the experiments deliver — the
          thesis pilots CA-style control, it does not prove emergence
          at scale.
\end{itemize}

% =============================================================
\section{Matter Does Matter}
\label{sec:matter-does-matter}
% =============================================================

\textbf{Role.} Establish embodiment / morphological computation
as the second conceptual pillar — the body is a computational
participant, not a passive platform. Grounds the
\emph{Soft vs.\ Rigid} dialectic from \S1.1 and the
configuration-as-computation argument central to the thesis.

\textbf{POV.} The interesting design move is not ``soft instead of
rigid'' but \emph{which computation we choose to do in matter and
which we leave to a controller}.

% -----------------------------------------
\subsection{The classical separation of body and controller}
\label{subsec:classical-separation}

\textbf{Key arguments.}
\begin{itemize}
    \item Industrial robotics treats the body as a precise mechanical
          platform. The controller specifies torques; the body executes.
    \item This separation has been productive but it makes any
          computation done by the body \emph{invisible}, and treats
          compliance, friction, and material variability as noise to
          be rejected.
    \item Re-thinking this is a precondition for taking soft and
          modular bodies seriously as design objects.
\end{itemize}

\textbf{References (from your .bib).}
\cite{rusDesignFabricationControl2015} (as the canonical statement
of the soft-robotics counterposition).

% -----------------------------------------
\subsection{Morphological computation: the body as participant}
\label{subsec:morph-comp}

\textbf{Key arguments.}
\begin{itemize}
    \item Pfeifer \& colleagues' framing: a well-designed body
          \emph{offloads} computation that a controller would
          otherwise have to specify.
    \item Three exemplars to anchor the principle: passive dynamic
          walking; the universal soft / jamming gripper; physical
          reservoir computing in compliant bodies.
    \item Acknowledge the careful critique: not every offload is
          ``computation proper'' --- much is morphology that
          \emph{facilitates} control or perception. Draw the
          distinction carefully so the thesis claim is the right
          shape.
\end{itemize}

\textbf{References (from your .bib).}
\cite{pfeiferHowBodyShapes2007},
\cite{pfeiferNewRoboticsDesign2005},
\cite{mullerWhatMorphologicalComputation2017},
\cite{wangEmbodyingPhysicalComputing2026},
\cite{aguilarReviewLocomotionRobophysics2016}.

\textbf{Consensus candidates worth adding.}
\begin{itemize}
    \item \emph{Hauser et al.\ 2023, IEEE Control Systems} ---
          ``Leveraging Morphological Computation for Controlling
          Soft Robots''. Translates the principle into a control-
          design vocabulary; useful framing if you want a control-
          aware reader to follow you.
    \item \emph{Hauser et al.\ 2024, Device} --- ``Morphological
          computation: past, present and future''. Recent (2024)
          short survey; good for an up-to-date anchor.
    \item \emph{Laschi et al.\ 2016, IEEE RAM} --- ``Lessons from
          animals and plants: the symbiosis of morphological
          computation and soft robotics''. Directly bridges the two
          literatures you are joining.
    \item \emph{Müller \& Hoffmann 2017} (you already have it) is the
          critical/cautious counterweight --- cite it precisely so
          you signal you are aware of the boundary between
          \emph{computation proper} and \emph{morphology that aids
          control}.
\end{itemize}

% -----------------------------------------
\subsection{Morphing matter and programmable self-assembly}
\label{subsec:morph-matter}

\textbf{Key arguments.}
\begin{itemize}
    \item A constructive sibling to morphological computation: matter
          itself can be designed to change form in response to stimuli
          or assembly rules.
    \item This pushes the design question from ``what compliance \emph{is}
          available?'' to ``what compliance do we \emph{author}?''
    \item Useful as bridging vocabulary: your inflatable modules are
          themselves a form of programmable matter at vertex scale.
\end{itemize}

\textbf{References (from your .bib).}
\cite{tibbitsThingsFallTogether2021},
\cite{cademartiriProgrammableSelfassembly2015},
\cite{wangRoboMatterReconfigurableMultifunctional2024},
\cite{adminNewProgrammableMaterials2022},
\cite{schaffner3DPrintingRobotic2018},
\cite{pattersonMethod3DPrinting2023}.

% -----------------------------------------
\subsection{Implication: the design question changes}
\label{subsec:design-question-changes}

\textbf{Key arguments.}
\begin{itemize}
    \item Reframe: instead of ``what controller do we need?'', ask
          ``what body would shrink the controller's task?''
    \item This reframe is the conceptual lens this thesis adopts.
    \item Important caveat: most demonstrations of morphological
          computation are on \emph{single-body, monolithic} systems.
          The question of how morphological computation operates in
          \emph{modular} bodies, where modules are coupled through
          their material, has received much less attention.
\end{itemize}

\textbf{Closing transition.} The principle of matter-as-computation
is well established in monolithic soft systems and compliant
single-body machines. What kinds of robots, then, instantiate this
principle in \emph{modular} form --- where the body is composed
of comparable units, and the configuration of those units is itself
a design variable? Two robotic traditions have approached this
question from opposite sides; \S2.3 reads each in turn, then
identifies the region they leave open.

\textbf{Watch out for.}
\begin{itemize}
    \item Risk of overlap with \S2.3. Keep this section
          \emph{conceptual}; save specific comparable systems
          (Wait, Chin, Nozaki, Usevitch, Yu) for the next section.
    \item If you borrow Tibbits' \emph{Robot without Robot} framing
          for \S2.3's title, cite him here once explicitly.
\end{itemize}

% =============================================================
\section{Robot Unlike Robot}
\label{sec:robot-unlike-robot}
% =============================================================

\textbf{Role.} \emph{Position} the thesis against existing robotic
systems. This is where the literature has to deliver the gap.
Title borrowed from Tibbits --- cite \cite{tibbitsThingsFallTogether2021}
in the section's first paragraph.

\textbf{POV.} Two strong but lopsided traditions: soft robotics
takes matter seriously and comparability lightly; modular
reconfigurable robotics does the inverse. The interesting
position is the one neither has fully claimed.

% -----------------------------------------
\subsection{Soft robotics --- compliance without comparability}
\label{subsec:soft-robotics}

\textbf{POV / closing point.} These systems show what compliance
\emph{can do}. They do not provide a way to compare configurations
under shared geometric constraints, because there \emph{are no
shared structural primitives} across designs.

\textbf{Key arguments.}
\begin{itemize}
    \item Soft robotics' design language is largely \emph{monolithic} ---
          a single continuous body whose behaviour is hard to compare
          across designs.
    \item Pneumatic / inflatable variants are the closest cousin of
          this thesis at the level of actuation, but stay within the
          monolithic frame.
    \item Authoring tools and fabrication methods (3D printing,
          inflatables-by-design) define the practical possibility
          space.
\end{itemize}

\textbf{References (from your .bib).}
Foundational review: \cite{rusDesignFabricationControl2015}.
Pneumatic/inflatable single bodies:
\cite{waitSelfLocomotionSpherical2010},
\cite{usevitchUntetheredIsoperimetricSoft2020},
\cite{kimSimpleTripodMobile2019},
\cite{boothSurfaceActuationSensing2021}.
Fluidic / self-synchronizing units:
\cite{mousaMultifunctionalFluidicUnits},
\cite{eguchiFluidicSwarmEmbodimentSwarm2025}.
Inflatable design / fabrication:
\cite{yang_snapinflatables_2024},
\cite{wangPneuFabDesigningLowCost2023a},
\cite{ouTEI2016Studio2016},
\cite{schaffner3DPrintingRobotic2018},
\cite{pattersonMethod3DPrinting2023}.

\textbf{Consensus candidates worth adding (optional).}
\begin{itemize}
    \item \emph{Shepherd et al.\ 2011, PNAS} --- ``Multigait soft robot.''
          The 1971-citation foundational pneumatic soft-robot paper.
          Even a passing cite signals you've placed yourself in the
          lineage.
    \item \emph{Hauser et al.\ 2023} (already suggested for \S2.2) ---
          alternatively could land here as ``soft + control.''
\end{itemize}

% -----------------------------------------
\subsection{Modular reconfigurable robotics --- comparability without compliance}
\label{subsec:mrr}

\textbf{POV / closing point.} These systems show what comparable
design spaces look like, and how powerfully they can be reasoned
about. They block exactly the passive material coupling that makes
soft systems interesting in the first place.

\textbf{Key arguments.}
\begin{itemize}
    \item MRR offers exactly what soft robotics lacks --- comparable
          units, design spaces that can be enumerated and studied,
          policy frameworks that generalize across designs.
    \item But almost all MRR is \emph{rigid}; reconfigurability lives
          in mechanical re-attachment or joint angles, not in the
          body's material behaviour.
    \item Recent learning-based work (modular RL, GAN-driven design
          search) has only deepened this rigid commitment because
          differentiable / RL pipelines presume well-modelled bodies.
\end{itemize}

\textbf{References (from your .bib).}
Reviews: \cite{ahmadzadehModularRoboticSystems2015},
\cite{liangModularReconfigurableRobots2026}.
Canonical rigid modular systems:
\cite{romanishinMBlocksThreeDimensional2018},
\cite{rubensteinProgrammableSelfassemblyThousandrobot2014},
\cite{tuFreeSNFreeformStrutnode2022},
\cite{yiReconfigurableRobotControl2023},
\cite{yiImprovingRobotCapabilities},
\cite{chinFunctionFollowsForm2023}.
Learning-based modular design:
\cite{whitmanLearningModularRobot2023},
\cite{huModularRobotDesign2022},
\cite{38942241ModularMobile},
\cite{AgileLeggedLocomotion}.
Distributed / bio-inspired modular control:
\cite{zhangBioInspiredDistributedNeural}.

\textbf{Consensus candidates worth adding.} None essential. Yim 2007
(modular self-reconfigurable robots review) is conventionally cited
here and isn't in your bib --- \emph{add only if your committee
expects the canonical Yim citation}; \cite{ahmadzadehModularRoboticSystems2015}
already covers the review function.

% -----------------------------------------
\subsection{At the boundary --- soft \emph{and} modular \emph{and} reconfigurable}
\label{subsec:boundary}

\textbf{POV.} A handful of systems sit at or near the intersection.
Each pays a different price; none asks the question this thesis
asks.

\textbf{Key arguments.}
\begin{itemize}
    \item Polyhedral-geometry-as-design-primitive is rare; closest
          neighbours are Mochibot (rigid, max-DOF) and tensegrity
          icosahedra (cabled, not modular in the swappable-unit sense).
    \item Particle robots are soft-edged and collective but their
          ``configuration'' is dynamic and emergent, not designed.
    \item Volume-changing modular systems --- AuxBots, electrohydraulic
          modules --- are reconfigurable and modular but rigid.
    \item The closest soft + modular systems with \emph{vertex-style}
          modular geometry are very recent and not yet positioned as
          a comparable design family.
\end{itemize}

\textbf{References (from your .bib).}
\cite{nozakiContinuousShapeChanging2018} (Mochibot),
\cite{chinFlipperStyleLocomotionStrong2023} (AuxBots),
\cite{liParticleRoboticsBased2019} (particle robots),
\cite{wangRoboMatterReconfigurableMultifunctional2024},
\cite{younPneuBotsModularInflatables2022} (PneuBots --- soft modular
inflatable, the nearest sibling in your bib),
\cite{yoderHexagonalElectrohydraulicModules2024},
\cite{boothSurfaceActuationSensing2021} (tensegrity skin),
\cite{tibbitsThingsFallTogether2021} (programmable matter framing),
\cite{suzukiShapeBotsShapechangingSwarm2019}.

\textbf{Consensus candidates worth adding (IMPORTANT).}
\begin{itemize}
    \item \textbf{Wharton et al.\ 2023, IEEE RA-L --- Tetraflex /
          Polyflex.} A \emph{tetrahedral pneumatic-bellows soft truss
          robot} that rolls, crawls, and bounds. This is a direct
          neighbour of your tetra/octa/cube comparison and is not in
          your bib. Strongly recommend adding --- it is either the
          nearest competitor or a missing comparator depending on
          how you frame the gap. Worth a paragraph of its own.
    \item \emph{Wang et al.\ 2022, T-MECH} --- ``Tetrahedral
          Soft-Legged Robot for Multigait Locomotion.'' Less directly
          comparable (legged not vertex-modular) but uses the
          tetrahedron framing.
    \item \emph{Robertson et al.\ 2020, IJRR} --- ``Pneumagami''
          modular origami pneumatic soft robot. Useful for the
          ``modular soft pneumatic'' lineage even though it is
          manipulator-style, not locomotor.
    \item \emph{Pagliocca et al.\ 2024} --- modular reconfigurable
          rotary soft pneumatic actuators. Demonstrates modular
          soft pneumatic units that can recombine.
\end{itemize}

\textbf{Comparison table (anchor of this subsection).} Lead with the
table, let the prose support it.

\begin{table}[h]
\centering
\small
\begin{tabular}{lcccc}
\toprule
System & Modular? & Soft? & Topological comparison? & Locomotion via inflation? \\
\midrule
Wait et al.\ 2010 & No & Yes (pneumatic) & No & Yes \\
Usevitch et al.\ 2020 & No & Yes & Single config & Partial \\
Chin et al.\ 2023 (AuxBot) & Yes & No (rigid auxetic) & Limited & No \\
Nozaki et al.\ 2018 (Mochibot) & No (monolithic) & Partial & No (single shape) & No \\
Yu et al.\ 2026 (MRR) & Yes & No & Yes (rigid) & No \\
Li et al.\ 2019 (particles) & Yes (collective) & Partial & Emergent only & No \\
Wharton et al.\ 2023 (Tetraflex) & Yes (truss) & Yes (bellows) & Single config & Yes \\
\textbf{This thesis} & \textbf{Yes (vertex modules)} & \textbf{Yes (pneumatic)} & \textbf{Yes (cospherical family)} & \textbf{Yes} \\
\bottomrule
\end{tabular}
\caption{Comparable systems against the thesis along four axes.}
\label{tab:comparison}
\end{table}

% -----------------------------------------
\subsection{The gap}
\label{subsec:gap}

\textbf{POV.} A precise, actionable gap statement that bridges
to Chapter~3.

\textbf{Synthesis (draft prose, can be promoted to body text).}
Existing soft robots demonstrate that matter performs computational
work, but their monolithic form makes it difficult to ask how the
\emph{arrangement} of compliant elements shapes that work. Existing
modular robots demonstrate that arrangement can be studied as a
design variable, but their rigid construction blocks the passive
material coupling that makes the question of arrangement interesting
in the first place. Between these two traditions lies a thinly
populated region: systems that are simultaneously soft, modular, and
reconfigurable, in which both compliance and configuration can be
varied and compared on shared terms.

It is in this region that this thesis is sited. The chapters that
follow develop a geometric design framework --- the \emph{cospherical
polyhedral primitive} --- that operationalizes this region, allowing
different topological configurations of identical inflatable modules
to be designed, fabricated, simulated, and tested under one
comparative language.

\textbf{Closing transition (to Ch.\ 3).} What this region requires
is a design framework that holds compliance and configuration
together --- that lets the geometric arrangement of soft modules be
specified and varied without losing comparability across designs.
Such a framework is the subject of the next chapter.

% =============================================================
\section*{Cross-section reminders}
\label{sec:reminders}
% =============================================================

\begin{itemize}
    \item \textbf{No double-citing the same paper for the same argument
          across sections.} Pfeifer \& Bongard 2007 belongs in \S2.2;
          if it appears in \S2.1 or \S2.3, it should be carrying a
          different argument.
    \item \textbf{Length budget:} \S2.1 $\approx$ 3--4 pages, \S2.2 $\approx$
          3--4 pages, \S2.3 $\approx$ 5--7 pages (longest, since it
          delivers the gap). Total chapter $\approx$ 12--15 pages.
    \item \textbf{Transition discipline:} every section ends by naming
          what the next section addresses; the chapter as a whole
          ends by naming what Ch.~3 does.
    \item \textbf{Dialectic callbacks (per your preference):} each
          section opener names the dialectic from \S1.1 it picks up.
\end{itemize}

% =============================================================
\section*{Open questions for us to settle before drafting}
\label{sec:open-questions}
% =============================================================

\begin{enumerate}
    \item \textbf{Tetraflex.} Add Wharton et al.\ 2023 to the bib?
          If yes: position as nearest competitor (it is a tetrahedral
          soft pneumatic truss --- the very configuration you found
          to fail) or as ``proves the boundary is being approached
          from another angle''?
    \item \textbf{Foundational classics.} Do you want Anderson 1972,
          Reynolds 1987, Wolfram 2002, McGeer 1990 in the bib? Your
          current set works without them; adding them costs little
          but expands the bibliography.
    \item \textbf{Yim 2007 modular review.} Conventionally cited in
          MRR sections; not in your bib. Your committee's expectation
          here?
    \item \textbf{Hauser 2023 / 2024 morphological computation
          updates.} Add to \S2.2 or skip? Recent surveys would
          modernize the section without much rewrite.
    \item \textbf{Cybernetics thread.} \cite{glanvilleSECONDORDERCYBERNETICS}
          and \cite{sheridanHumanRobotInteractionStatus2016} are in the
          bib but probably belong in the reflection chapter, not here.
          Confirm we just acknowledge in the chapter intro?
    \item \textbf{Comparison table placement.} Currently in
          \S\ref{subsec:boundary}. Could move to \S\ref{subsec:gap}
          or to a chapter-opener visual. Your call.
\end{enumerate}
